\def\velkostpisma{9.9}
\def\velkostriadku{12.5}
\input makra.tex % nacitanie Ivanom pripravenych nastaveni a prikazov
\hyphenation{star-šov-stvo} % rozdelenie slov na konci riadku, treba tu uviest slova, ktore sam nepozna

\spravodaj{11}{2026}


\clanok{{\it ČO SA DEJE V~NOVEMBRI}}


\clanok{Dve stretnutia sestier -- 4.~11. a 18.~11.}
Sestry sa stretnú v~novembri dvakrát. Prvé stretnutie sa uskutoční v~stredu 4.~11. o~17.30~hod. na Zrínskeho~2, jeho témou bude Svetový deň modlitieb baptistických žien. Druhé stretnutie sa uskutoční v~stredu 18.~11. o~17.30~hod. v~budove SED na Partizánskej~2. Našou hosťkou bude ev.~farárka, s.~Anna Činčuráková. Jej tému doplníme čoskoro. Sestry všetkých generácií, ste srdečne vítané.


\clanok{Rozhovory na Palisádach -- 7.~11.}
Starnúť a nestratiť chuť do života patrí medzi najkrajšie, ale aj najťažšie životné úlohy. Nie preto, že by sme mali byť stále veselí, mladí a plní energie, ale preto, aby sme aj vtedy, keď sa mení naše telo, ubúdajú nám sily alebo sa niektorí ľudia z~nášho života vytrácajú, stále cítili, že stojí za to žiť. O~tom, ako sa to všetko odohráva v~živote starších ľudí, porozpráva v~sobotu 7.~11. o~9.30~hod. Bea Dobová, riaditeľka Zariadenia pre seniorov SED v~Bratislave.


\clanok{Spoločný obed v~SED -- 8.~11.}
Milí bratia a sestry, priatelia, po letnej a septembrovej prestávke vás opäť pozývame na spoločný obed. Pre jesenné prázdniny ho presúvame z~prvej novembrovej nedele na druhú, teda na 8.~11. o~12.00~hod. (môžete prichádzať hneď po zhromaždení). Ak máte chuť obedovať a rozprávať sa, príďte so svojím obedom a vezmite so sebou rodinu alebo priateľov. Vždy sme pripravení nasýtiť aj tých, ktorí by v~ten deň obed nemali a chceli by sa k~nám pridať.
\autor{s~láskou „obedový tím“}


\clanok{Adventné bohoslužby vo štvrtky na Palisádach -- od 26.~11.}
Podobne ako minulý rok plánujeme vo štvrtky (namiesto biblickej hodiny) adventné bohoslužby v~našej modlitebni na Palisádach. Začínať sa budú o~18.00~hod. a viesť ich bude br.~kazateľ Dávid M.~Chuchút. Uskutočnia sa 26.~11., 3.~12., 10.~12. a 17.~12. Vítaní sú všetci, ktorí chcú
v~tichu aj pri hudbe rozjímať nad Božím Slovom a Božím dielom v~našich životoch.


\clanok{Matuzalem revival -- 28.~11.}
Srdečne vás pozývame na koncert Matuzalem revival, ktorý sa uskutoční pri príležitosti životného jubilea nášho br.~Sláva Kráľa v~sobotu 28.~11. o~17.00~hod. na Palisádach.
Príďte sa potešiť a potešiť svojou prítomnosťou aj brata Sláva.


\clanok{{\it SPRÁVY Z~OKTÓBRA}}


\clanok{Správy zo staršovstva}
Drahí bratia a sestry,
staršovstvo zboru BJB Palisády sa v~októbri stretlo dvakrát, a to 13.~10. a 27.~10. Na prvom stretnutí sme sa v~prvej časti vyučovali a diskutovali o~knihe Starší zboru. V~druhej časti sme sa venovali pastorácii. V~tretej časti sme sa zaoberali agendou zboru: kontrolou uznesení, novými informáciami o~budúcnosti F.~Barkócziho, zhodnotením vďakyvzdania v~SÚZA, zhodnotením KDZ a prípravou ZČZ. V~rámci bodu Rôzne sme hovorili o~kázňach od Adventu po Veľkú noc, Rozhovoroch na Palisádach, spoločnom obede v~SED na Partizánskej a iniciatíve D.~Uhrina „Slovensko a zmierenie“.

Na druhom stretnutí staršovstva sme sa v~prvej časti vyučovali a diskutovali o~knihe Starší zboru. V~druhej časti sme sa venovali pastorácii. V~tretej časti sme sa zaoberali agendou zboru: kontrolou uznesení, vyhodnotením ZČZ a KDZ, dotazníkom a programom zhromaždení. V~rámci bodu Rôzne sme hovorili o~konferencii CBS, zbierke pre Dobrodinu a rekonštrukcii modlitebne. Najbližšie stretnutie staršovstva bude 10.~11.
\autor{Marcel Maďar}


\clanok{Oblastné vďakyvzdanie -- nedeľa 11.~10.}
Ďakujeme za každú ruku priloženú k~dielu. Za počúvanie o~vašej vďačnosti, za spievanie piesní chvál a vďaky, za vzájomné spoločenstvo, za deti, za koláčiky, úsmevy či rozhovory.
\vfill\break


\clanok{Rozhovory na Palisádach}
Po letnej prestávke sme pokračovali v~sérii pravidelných podujatí „Rozhovory na Palisádach“ s~hosťkou MUDr.~Andreou Škripekovou.
Andrea predstavila medicínsky odbor paliatívnej starostlivosti, ktorého cieľom je, aby človek žil čo najlepšie aj s~vážnou chorobou.
Paliatívna medicína lieči, čo človeka trápi. Dnes už dokáže poskytnúť úľavu od bolesti a únavy, ale aj od strachu a neistoty, ktoré s~chorobou prichádzajú.
Pacienti paliatívnej starostlivosti majú možnosť rozhodovať sa a je im umožnené prepájať medicínske možnosti s~hodnotami človeka.
Naša hosťka v~diskusii reagovala na mnoho otázok, vrátane tých, ktoré sa venovali eutanázii, pretože paliatívna starostlivosť ponúka riešenia aj pre ľudí, ktorí premýšľajú týmto spôsobom.
Moderná paliatívna starostlivosť tiež dokáže reagovať na potrebu bezpečia a dôstojnosti svojich pacientov a mnohí dokážu s~touto podporou plnohodnotne žiť aj niekoľko rokov.
\autor{Martin Simon}


\clanok{Konferencia EBWU (European Baptist Women United) v~Prahe}
Únia baptistických žien Európy spája baptistické ženy z~52~štátov Európy a Blízkeho východu. Tohto roku, v~dňoch 8.~--~11.~10., mali české a slovenské sestry vzácnu príležitosť hostiť konferenciu EBWU v~Prahe. Táto konferencia sa koná každý rok v~inej krajine Európy alebo Blízkeho východu.
Témou konferencie bolo „NEOTRASITEĽNÁ“ podľa žalmu: „Nikdy sa nezachveje. Vo večnej pamäti bude spravodlivý. Nezľakne sa zlej správy, jeho srdce je pevné, dúfa v~Hospodina.“ (Ž~112,6-7)

Hlavnou rečníčkou bola Dr.~Lina AbuJamra, ktorá na príbehu Mojžišovho života ukázala, že skutočná viera sa prejavuje vernosťou Bohu aj uprostred neistoty, bolesti a čakania.
Mojžišov život nebol bez zlyhaní, strachu či sklamaní. Štyridsať rokov strávených na púšti sa však preňho stalo časom Božej prípravy. Rečníčka nám pripomenula, že núdza je barometrom našej viery, no žiadna núdza nie je taká veľká, aby ju Boh nevidel alebo aby nedokázal naplniť naše potreby.
Dotkla sa aj bolesti, ktorú prinášajú medziľudské vzťahy, dokonca aj v~cirkvi. Povzbudila nás, aby sklamania neotriasli našou vierou, pretože Boh zostáva naším verným priateľom aj vtedy, keď sa cítime osamelí a nepochopení.
\autor{Jarka Cihová}
\vfill\break


\clanok{Októbrové Dobronovinky}
Ako sme už avizovali, mesačné informácie o~činnosti Dobrodiny budeme pravidelne prinášať v~Spravodaji. Prečítať si ich môžete aj prostredníctvom odkazu v~e-mailových oznamoch.
Správy budú vždy s~mesačným oneskorením. Nech sa páči, tu sú októbrové Dobronovinky:
\vskip0.01em
{\em Keď Boh rozmnožuje, čo Mu zveríme.}
\vskip0.01em
Milí bratia a sestry, dnes sa s~vami chceme podeliť o~jeden príbeh z~úplných začiatkov Dobrodiny.

V~roku~2022 nám istý Bohu známy darca zo Slovenska poslal na účet 30~000~€ na podporu Dobrodiny. Tento dar (spolu s~ďalšími darmi, za ktoré sme takisto vďační) pomohol rozbehnúť celý projekt. Predovšetkým bol však pre nás uistením, že nám niekto dôveruje v~úplných začiatkoch až tak, že je ochotný venovať veľkú sumu peňazí na toto začínajúce Božie dielo.

Minulý týždeň sme podpísali s~Ministerstvom práce, sociálnych vecí a rodiny zmluvu o~nenávratnom finančnom príspevku na vybudovanie špecializovaného zariadenia. Príspevok je vo výške viac ako 3~mil.~€, teda viac ako stonásobok spomínaného daru.
Vďaka tomuto príspevku môžeme pokračovať v~ďalšej etape a smerovať k~tomu, aby sme v~roku~2028 otvorili špecializované zariadenie pre ľudí, ktorí budú potrebovať našu pomoc a starostlivosť.

V~Dobrodine už niekoľko rokov zažívame, že keď Pánu Bohu ponúkneme to, čo máme, On to vie rozmnožiť mnohonásobne. A~tak aj Vás pozývame, aby ste sa pridali a podporili toto dielo čiastkou (či už jednorázovou alebo pravidelnou mesačnou platbou), ktorú Vám Pán Boh položí na srdce, s~vierou, že On ju dokáže rozmnožiť tak, aby sme o~niekoľko rokov mohli dokončiť všetky stavby Dobrodiny a otvoriť dvere ľuďom, ktorí budú naše služby potrebovať.

Bratislava UŽ spieva gospel a predaj lístkov sa spustil!
Bratislava spieva gospel má za sebou prvý spoločný nácvik -- a hneď s~rekordnou účasťou v~histórii BSG! V~sále sa zhromaždilo 250~spevákov a hneď počas prvého večera spolu nacvičili 4~piesne. Ďalších osem spoločných nácvikov ich ešte čaká.
Ako bude niekoľko stoviek hlasov znieť spolu, si môžete prísť vypočuť
29.~11. v~PEUGEOT aréne. A~my budeme radi, keď sa tam stretneme aj s~vami, pretože tento rok bude súčasťou príbehu BSG aj Dobrodina. Lístky sú už v~predaji na \ulink[https://www.ticketportal.sk/]{www.ticketportal.sk}.
Ako nás môžete podporiť?
Vaša podpora môže mať mnoho podôb -- modlitbu, finančný dar, zdieľanie Dobrodiny, či váš čas a schopnosti. Každú z~nich si úprimne vážime.
Po novom nás môžete podporiť aj cez platformu \ulink[https://www.darujme.sk/]{Darujme.sk}.
„Vždy, keď sa modlíme, ďakujeme za vás Bohu, Otcovi nášho Pána Ježiša Krista. Veď sme počuli o~vašej viere v~Ježiša Krista a o~láske, ktorú máte k~všetkým svätým pre nádej, ktorá je vám pripravená v~nebesiach a o~ktorej ste už skôr počuli v~pravdivom slove evanjelia, ktoré k~vám prišlo. Rovnako ako na celom svete aj u~vás prináša ovocie a rastie od toho dňa, keď ste počuli a poznali Božiu milosť v~pravde.“
(Kol~1,3-6)
\autor{s~vďačnosťou, tím Dobrodiny}


\clanok{Zbierky v~októbri}
\table{lr}{
Na misiu           &   0~€\cr
Na investičný fond & 575~€\cr
Na Dobrodinu       & 700~€\cr
}
\vskip1em

Pripomíname, že zbierky sa u~nás konajú pravidelne takto:
\vskip-1ex\begitems
* každú 2.~nedeľu v~mesiaci je zbierka venovaná misii a
* každú 4.~nedeľu je zbierka na investície.
\enditems

Prispieť je možné formou zbierky počas 4.~nedieľ v~mesiaci alebo bankovým prevodom na zborový účet, prípadne aj osobne s.~Ľubke Kohútovej.

Bankové spojenie: SK36 0900 0000 0000 1147 1836, SWIFT: GIBASKBX




\n 5.	11.	Katarína	VALENTOVÁ;
\n 6.	11.	Elena	PRIBULOVÁ;
\n 6.	11.	Eva	SYČOVÁ;
\n 9.	11.	Radovan	PAULEN;
\n 15.	11.	Bohumila	ŠALINGOVÁ;
\n 19.	11.	Dávid	PRIBULA;
\n 22.	11.	Peter	PRIBULA;
\n 25.	11.	Petra	ŠALINGOVÁ -- FOSSETI;
\n 27.	11.	Judita	KOLÁŘIKOVÁ;
\n 29.	11.	Jaroslav	KRÁĽ;
\narodeniny


\program{
\p  1 ; ne ;  9.30 ; Bohoslužby (J.~Szőllős + VP) ;.;;
\p  2 ; po ; 18.00 ; Modlitebná skupinka ;.;;
\p  3 ; ut ; 15.00 ; Biblická hodina pre seniorov (P.~Pivka) ;.;;
\p  4 ; st ; 17.30 ; Stretnutie sestier ;.; English Bible study ;
\p  5 ; št ; 18.00 ; Biblická hodina (J.~Szőllős) ;.;;
\p  6 ; pi ; 17.30 ; Dorast ;.;;
\p  7 ; so ;  9.30 ; Rozhovory na Palisádach ; 18.00 ; Mládež ;
\p  8 ; ne ;  9.30 ; Bohoslužby (D.~M.~Chuchút) ; 12.00 ; Spoločný obed v~SED ;
\p  9 ; po ; 18.00 ; Modlitebná skupinka / skupinka „Rastieme v~Kristovi“ ;.;;
\p 10 ; ut ; 15.00 ; Biblická hodina pre seniorov (P.~Pivka) ;.;;
\p 11 ; st ;.;;.;;
\p 12 ; št ; 18.00 ; Biblická hodina (D.~M.~Chuchút) ;.;;
\p 13 ; pi ; 17.30 ; Dorast ;.;;
\p 14 ; so ; 18.00 ; Mládež ;.;;
\p 15 ; ne ;  9.30 ; Bohoslužby (F.~Barkóczi) ;.;;
\p 16 ; po ; 18.00 ; Modlitebná skupinka ;.;;
\p 17 ; ut ; 15.00 ; Biblická hodina pre seniorov (P.~Pivka) ;.;;
\p 18 ; st ; 17.30 ; Stretnutie sestier v~SED ;.;;
\p 19 ; št ; 18.00 ; Biblická hodina (D.~M.~Chuchút) ;.;;
\p 20 ; pi ; 17.30 ; Dorast ;.;;
\p 21 ; so ; 18.00 ; Mládež ;.;;
\p 22 ; ne ;  9.30 ; Bohoslužby (J.~Szőllős) ;.;;
\p 23 ; po ; 18.00 ; Modlitebná skupinka / skupinka „Rastieme v~Kristovi“ ;.;;
\p 24 ; ut ; 15.00 ; Biblická hodina pre seniorov (P.~Pivka) ;.;;
\p 25 ; st ;.;;.;;
\p 26 ; št ; 18.00 ; Adventné bohoslužby na Palisádach (D.~M.~Chuchút) ;.;;
\p 27 ; pi ; 17.30 ; Dorast ;.;;
\p 28 ; so ; 17.00 ; Matuzalem revival ; 18.00 ; Mládež ;
\p 29 ; ne ;  9.30 ; Bohoslužby (D.~M.~Chuchút) ;.;;
\p 30 ; po ; 18.00 ; Modlitebná skupinka ;.;;
}
\tiraz
\bye
